
\section{Additional Experiments}

\noindent{\textbf{Relative Pose Estimation.}}
Following prior work~\cite{sarlin2020superglue, edstedt2024roma}, 
we additionally evaluate the quality of the predicted camera poses using the relative rotation accuracy metric. Specifically, we compare DUSt3R~\cite{wang2024dust3r}, MASt3R~\cite{leroy2024mast3r}, and VGGT~\cite{wang2025vggt} against AnySplat and \OURMODEL, which are trained with multi-view supervision on RE10K.
We construct the evaluation sets on RE10K and ACID following the same protocol as NoPoSplat~\cite{ye2024nopo}, but, unlike NoPoSplat, we directly use the camera parameters predicted by our model without any test-time optimization. 
As shown in \cref{sup_tab:camera_pose}, our model, trained on the RE10K dataset, achieves higher accuracy across all thresholds. 
Furthermore, despite not being trained on the ACID dataset~\cite{liu2021acid}, \OURMODEL\, generalizes effectively to this unseen dataset, achieving the best performance overall while outperforming all baselines at $10\degree$ and $20\degree$ and remaining competitive at $5\degree$. 
These results indicate that, although our geometry backbone is initialized from VGGT, the proposed architecture and training scheme further refine and extend its geometric reasoning, leading to consistently stronger pose estimation performance.

\noindent{\textbf{Computational cost of Spatially Adaptive Gaussian Allocation.}}
\cref{sup_tab:adaptive_cost} reports the computational overhead of introducing Spatially Adaptive Gaussian Allocation. The peak allocated VRAM increases by only $1.8\%$, and the inference time increases by $10.1\%$. As discussed in the \textit{Multi-Scale Prediction} paragraph of \cref{sec3.2:SADA} in the main paper, the proposed design predicts multi-scale maps efficiently, making adaptive Gaussian allocation possible with minimal extra cost.
